\subsection{運動の記述}
    1次元トラップポテンシャル中のイオンの位置および速度の時間発展を、古典力学に基づく運動方程式によって記述する。
    本研究では、イオンの重心運動を量子化せず、位置 $x(t)$ および速度 $v(t)$ を連続的な実数値として扱う半古典近似を用いる。
    したがって、運動状態は Schr\"odinger 方程式や振動量子数（phonon）ではなく、ニュートン型の運動方程式に従って時間発展する。
\subsubsection{ポテンシャル束縛}
    イオンを任意の領域にとどめるために、調和振動子ポテンシャルをかける。
    ポテンシャル中心を原点と取ることで、各粒子に以下の力$F_1$が加わるといえる。
    \begin{equation}
        F_{1n} = -k_nx_n
    \end{equation}
    ここでk_nはトラップ周波数$f_{trap}$,粒子質量$m_n$に対して
    \begin{equation}
        k_n = (2\pi f_{trap})^2 m_n
    \end{equation}
    とする。\\
\subsubsection{ローレンツ反発}
    イオン同士はクーロン相互作用によって力を及ぼしあう。同種のイオンを想定する場合、イオンは同じ電荷を持つため反発力が生まれる。\\
    各粒子に加わるクーロン相互作用の力は素の粒子を除く全ての粒子との相互作用による力の和であるから、合力は粒子数をMとして次の$F_2$であるといえる。\\
    \begin{equation}
        F_{2n} = \sum_{l=0}^M a_{ln} \frac{\alpha}{(x_l-x_n)^2}
    \end{equation}
    ここで$i<j$のとき$x_i < x_j$であるとすると、$a_{ln}$は次のように定義される。
    \begin{equation}
        a_{ln} =
        \begin{cases}
            1 & \text{if } l < n \\
            0 & \text{if } l = n \\
            -1 & \text{if } l > n
        \end{cases}
    \end{equation}
    \begin{equation}
        \alpha=\frac{q_l q_n}{4\pi \epsilon}
    \end{equation}
    特に真空中の一価のイオンの集合を考えた場合$\alpha=2.3\times10^-28N/m^2$である。
\subsubsection{レーザー冷却}
     ドップラー冷却による平均の力を考える。\\
     飽和パラメータ$S_0$, 冷却レーザー周波数$\gamma$, デチューニング$\delta $, レーザーの進行方向$\vec{k_l}$として、平均の力$<F_3>$は次のようにかける。
    \begin{equation}
        <F_3> = \hbar\gamma\frac{\frac{S_0}{2}}{S_0+1+\frac{4}{\gamma}(\delta-\vec{k_l}\cdot\vec{v})^2}\cdot|\vec{k_l}|
    \end{equation}
    一次元中の系であるから内積は単に積で計算をした。
\subsubsection{反跳加熱}
     原子は励起時に放出する光子から運動量キックを受ける。放出方向は規定できないため、ランダムな方法に運動量を受けるとすると、この運動量は次のように記述できる。
    \begin{equation}
        p_{kick}=|p_{kick}|\cdot\cos\theta, (0\leq\theta\2\pi)
    \end{equation}
    励起が$\Delta t$sに一度起こるとすると、運動量の大きさは次のように記述できる。
    \begin{equation}
        |p_{kick}|=\hbar k_l\sqrt{\gamma\frac{\frac{S_0}{2}}{S_0+1+\frac{4}{\gamma}(\delta-\vec{k_l}\cdot\vec{v})^2}(\Delta t)}
    \end{equation}
    よって、$\Delta t$秒ごとに以下の運動量キックを加えればいい。
    \begin{equation}
        p_{kick}=\hbar k_l\sqrt{\gamma\frac{\frac{S_0}{2}}{S_0+1+\frac{4}{\gamma}(\delta-\vec{k_l}\cdot\vec{v})^2}(\Delta t)}\cdot\cos{\theta}
    \end{equation}
\subsubsection{運動方程式と加熱計算}
     以上の力から、運動方程式は次のように記述できる。
    \begin{multline}
        m_n\frac{d^2}{dt^2}x_n =F_1+F_2+<F_3> \\=-k_nx_n+\sum_{l=0}^M a_{ln} \frac{\alpha}{(x_l-x_n)^2}+\hbar\gamma\frac{\frac{S_0}{2}}{S_0+1+\frac{4}{\gamma}(\delta-\vec{k_l}\cdot\vec{v})^2}\cdot|\vec{k_l}|
    \end{multline}
    また、この運動方程式から計算した$\Delta t$秒における速度$v_{m,\Delta t,cal}$に対し、$\Delta t$秒ごとに加熱すると以下の式が成り立つ。
    \begin{equation}
        v_{n,\Delta t,new}=v_{n\Delta t,cal}+\frac{\hbar k_l}{m_n}\sqrt{\gamma\frac{\frac{S_0}{2}}{S_0+1+\frac{4}{\gamma}(\delta-\vec{k_l}\cdot\vec{v})^2}(\Delta t)}\cdot\cos{\theta}
    \end{equation}
    この運動方程式及び加熱計算を微小時間ごとに規定回数計算することで各粒子の運動を計算する。